En este artículo se comparte el desarrollo de Codexy, una plataforma innovadora diseñada para resolver el persistente problema de gestión de inventarios manual en grandes instituciones educativas y corporativas. En un mundo donde la transformación digital se ha convertido en una necesidad imperiosa, muchas organizaciones aún operan con sistemas obsoletos que generan pérdidas significativas debido a la falta de trazabilidad precisa y los inevitables errores humanos en el manejo de activos. Codexy surge como una solución integral que combina tecnologías modernas con prácticas de desarrollo eficientes, demostrando cómo la ingeniería de software puede transformar procesos operativos tradicionales.

El proyecto se fundamenta en metodologías Ágiles, permitiendo una adaptación rápida a los requerimientos cambiantes de los usuarios finales. La arquitectura elegida, basada en .NET 9 con Clean Architecture, busca un equilibrio óptimo entre robustez técnica y facilidad de mantenimiento, evitando la complejidad innecesaria de arquitecturas distribuidas como los microservicios. En su lugar, se implementó un sistema monolítico modular que facilita el despliegue y la escalabilidad inicial, complementado por una aplicación móvil híbrida desarrollada en Ionic que garantiza sincronización en tiempo real a través de SignalR.

Los resultados obtenidos evidencian que la aplicación de patrones de diseño establecidos, principios SOLID y medidas de seguridad integradas desde las primeras etapas del desarrollo no solo mejoran la eficiencia operativa actual, sino que también preparan el terreno para futuras integraciones con inteligencia artificial y tecnologías sostenibles. Este enfoque holístico no solo aborda los desafíos técnicos, sino que también considera aspectos culturales y organizacionales, asegurando una adopción efectiva por parte de los usuarios. Finalmente, se concluye que incluso el código más sofisticado resulta inútil si no se garantiza una implementación que considere las necesidades humanas y operativas de la institución, destacando la importancia de un enfoque centrado en el usuario en el desarrollo de software empresarial.